\documentclass[12pt]{article}
\usepackage{fullpage}
\usepackage[T1]{fontenc}
\usepackage[utf8]{inputenc}
\usepackage{lmodern}
\usepackage{microtype}
\usepackage{amsmath,amssymb}
\usepackage{hyperref}
\usepackage{enumitem}

\title{Descriptive Language, Realism, and the Case for Theoreticism}
\author{}
\date{}

\begin{document}
\maketitle

\begin{abstract}
This essay analyses the function of descriptive language regarded as an
instrument for talking about the world, and on that basis evaluates the
thesis that operationalism and instrumentalism should be replaced by
\emph{theoreticism}---the view that inquiry proceeds within a framework of
theories whose aim is genuine explanation, not the mere correlation of
observables. I argue for three connected conclusions. (i)~The
descriptive use of language carries, built into its very structure, a
commitment to a mind-independent state of affairs that statements aim to
fit; this is what distinguishes description from other uses of language
and is the precondition of the notions of truth and falsity on which
science depends. (ii)~This structural feature supplies a genuine, though
defeasible, argument for realism: to use language descriptively at all is
to presuppose that there is a way the world is, independently of our
descriptions, against which descriptions can succeed or fail.
(iii)~Realism so motivated undermines operationalism and instrumentalism
and supports theoreticism, because both rejected doctrines must either
covertly help themselves to the descriptive/realist resources they deny,
or else collapse the distinction between getting things right and merely
predicting, a distinction that the descriptive function of language makes
unavoidable. I am careful to mark where the argument is conclusive and
where it is only abductive, so that the strength of the final
endorsement of theoreticism is not overstated.
\end{abstract}

\section{Aim, method, and standard of success}

The problem asks me to do four things: analyse the descriptive function
of language; ask whether that function argues for realism; ask whether
such realism justifies replacing operationalism and instrumentalism by
theoreticism; and reach a determinate verdict. Because this is a
philosophical thesis and not a formal theorem, ``proof'' here means a
\emph{valid argument from premises that are either analytic of the
relevant concepts or else more plausible than their denials}, together
with an explicit accounting of where the argument is deductive and where
it is merely abductive. I will accordingly state premises explicitly,
draw the inference, and then test each premise against the strongest
objection I can find. A claim is allowed to stand only if it survives
that test; where it survives only as the best available explanation
rather than as a necessary truth, I say so.

I fix terminology first, because the entire argument turns on not letting
the key terms drift.

\begin{description}[leftmargin=1.6em,style=nextline]
\item[Descriptive (representative) use of language.] The use of a
declarative sentence to \emph{say how things stand}---to represent a
state of affairs as obtaining---so that what is said is correct if things
are as represented and incorrect otherwise.
\item[Realism.] The conjunction of two theses: (a)~\emph{ontological
realism}: there is a world whose constitution does not depend on our
beliefs, theories, or descriptions of it; (b)~\emph{semantic realism}:
the truth-value of a description is fixed by how that mind-independent
world is, not by our evidence, verification procedures, or predictive
convenience.
\item[Operationalism.] The doctrine that the meaning of a scientific
concept is \emph{exhausted by} the operations used to measure it: a
concept just \emph{is} the corresponding set of measuring operations.
\item[Instrumentalism.] The doctrine that scientific theories are not
candidates for truth or falsity but only \emph{instruments} (calculating
devices, inference tickets) for systematising and predicting observable
phenomena.
\item[Theoreticism.] The doctrine that all inquiry---including
observation and measurement---is conducted from within a framework of
theories, and that the proper aim of such theories is to \emph{explain},
i.e.\ to say truly what underlies and produces the phenomena, not merely
to correlate observables.
\end{description}

\section{The function of descriptive language as an instrument}

\subsection{The hierarchy of functions of language}

Language performs several distinguishable functions. Following the
analysis associated with B\"uhler and refined by Popper, we may
distinguish at least four, ordered by inclusion:

\begin{enumerate}[label=(F\arabic*),leftmargin=2.4em]
\item\label{F1} the \emph{expressive} (symptomatic) function: a verbal
utterance expresses or symptomises the inner state of the speaker;
\item\label{F2} the \emph{signalling} (release) function: an utterance
stimulates or releases a reaction in a hearer;
\item\label{F3} the \emph{descriptive} (representative) function: an
utterance describes a state of affairs, and is thereby true or false;
\item\label{F4} the \emph{argumentative} function: utterances stand in
relations of support, contradiction, and inference to one another.
\end{enumerate}

The ordering is one of inclusion in the following precise sense: whenever
a function higher on the list is present, the lower ones are present too,
but not conversely. \ref{F1} and \ref{F2} are shared with animal
signalling: a cry of fear expresses an inner state \ref{F1} and may warn
the herd \ref{F2} without describing anything. The decisive break is at
\ref{F3}. Only when an utterance \emph{describes}---purports to say how
things are---does the question of \emph{truth} arise; and only where
truth is in play can the \emph{argumentative} function \ref{F4} (which
trades in truth-preservation) get a grip. The higher functions are not
detachable embellishments: they are what make the lower functions
\emph{usable for inquiry at all}.

\begin{description}[leftmargin=1.6em,style=nextline]
\item[Lemma 1 (Truth presupposes description).] The predicates ``true''
and ``false'' apply to a use of language only insofar as that use is
descriptive in sense \ref{F3}.
\end{description}

\noindent\emph{Argument.} An utterance is true when things are as it says
they are, and false otherwise (Tarski's adequacy condition: the sentence
``snow is white'' is true iff snow is white). This biconditional has
application only where there is something the utterance \emph{says is the
case}---a represented state of affairs that may or may not obtain. A
purely expressive cry \ref{F1} or a purely releasing signal \ref{F2} says
nothing is the case; there is no represented state of affairs to be
compared with the world; so ``true'' and ``false'' lack purchase. Hence
truth-evaluability and descriptive use stand or fall together.
$\qquad\blacksquare$

\subsection{The internal structure of a description}

Lemma~1 lets us read off the internal structure that any descriptive use
must have. To describe is to deploy a representation $R$ that
\emph{purports to correspond} to a state of affairs $S$, in such a way
that the description \emph{succeeds} if $S$ obtains and \emph{fails} if it
does not. Three components are therefore constitutive of the descriptive
function:

\begin{enumerate}[label=(\roman*),leftmargin=2.2em]
\item a \emph{representation} $R$ (the content of what is said);
\item a \emph{represented item} $S$ (a state of affairs in the world);
\item a \emph{success condition}: $R$ is correct iff $S$ obtains, where
``$S$ obtains'' is not the same fact as ``$R$ is asserted'' or ``$R$ is
accepted''.
\end{enumerate}

\begin{description}[leftmargin=1.6em,style=nextline]
\item[Lemma 2 (Bipolarity / the possibility of error).] Any genuinely
descriptive use of language admits both correctness and error: it is
constitutive of $R$'s being a description that it could be \emph{false}.
\end{description}

\noindent\emph{Argument.} Suppose, for contradiction, that some use $R$
were descriptive yet could not be false---that nothing would count as
$R$'s misrepresenting how things stand. Then there would be no contrast
between $R$'s success condition obtaining and its failing to obtain; the
correctness of $R$ would not depend on anything beyond the making of $R$
itself. But then $R$ would not \emph{say how things stand} as opposed to
merely \emph{occurring}; it would have collapsed into the expressive or
signalling functions \ref{F1}--\ref{F2}, contradicting the assumption
that it is descriptive in sense \ref{F3}. Hence a descriptive use must be
\emph{bipolar}: capable of being right and capable of being wrong.
$\qquad\blacksquare$

Lemma~2 is the crux of the whole essay, so it is worth stating its force
exactly. It does not yet say that any particular description \emph{is}
true, nor that we can ever \emph{know} which descriptions are true. It
says something weaker but decisive: the very intelligibility of a
descriptive use of language requires a \emph{standard of correctness that
is not guaranteed by the act of describing}. There must be something the
description answers to.

\section{From the descriptive function to realism}

\subsection{The argument}

I now extract the argument for realism. Let $R$ be any descriptive use of
language (e.g.\ a measurement report, an observation statement, a
theoretical claim).

\begin{enumerate}[label=(P\arabic*),leftmargin=2.4em]
\item\label{P1} If $R$ is descriptive, then $R$ has a success condition
that is independent of the act of asserting or accepting $R$ (Lemmas~1
and~2).
\item\label{P2} A success condition independent of our assertion and
acceptance is a condition fixed by \emph{how things stand}, i.e.\ by some
state of affairs $S$ whose obtaining is not constituted by our describing
it.
\item\label{P3} The totality of such mind-independent states of affairs
that fix the success conditions of our descriptions is what ``the world''
denotes in semantic realism; and the thesis that descriptions answer to
it is semantic realism (Definition, \S1).
\item\label{C1} Therefore, to use language descriptively at all is to
treat the world as mind-independent in the sense of semantic realism: the
descriptive function \emph{presupposes} realism.
\end{enumerate}

The inference from \ref{P1}--\ref{P3} to \ref{C1} is valid: it is a chain
of biconditional unpackings of the term ``descriptive''. The premises
\ref{P1} and \ref{P3} are analytic of the notions of description and of
semantic realism respectively. The only premise that does substantive
work is \ref{P2}, and it requires defence against the obvious
anti-realist rejoinder.

\subsection{Defending \ref{P2}: ruling out the verificationist reduction
of the success condition}

The anti-realist (verificationist, operationalist, or instrumentalist)
will resist \ref{P2} by proposing to \emph{reduce} the success condition
of $R$ to something epistemic: $R$ is ``correct'' just when $R$ is (or
would be) verified, confirmed, warrantedly assertible, or predictively
useful. If that reduction succeeds, the success condition is not
independent of us, and \ref{C1} fails. So everything turns on whether the
success condition of a description can be identified with its
verification condition.

\begin{description}[leftmargin=1.6em,style=nextline]
\item[Lemma 3 (Non-identity of truth and verification).] The success
condition of a descriptive use of language cannot be identified with any
condition stated purely in terms of our verifying, confirming, or
predictively exploiting it.
\end{description}

\noindent\emph{Argument.} Consider any candidate epistemic surrogate
$V(R)$ (``$R$ is verified'', ``$R$ would survive our tests'', ``$R$ is
predictively reliable''). The statement $V(R)$ is itself a description: it
purports to say how things stand with our evidence and procedures, and by
Lemma~2 it too is bipolar---it can be \emph{mistaken}. We can
intelligibly say ``$R$ has been verified, yet $R$ is false'' and ``$R$ has
not been verified, yet $R$ is true''. These conjunctions are not
self-contradictory: the history of science is full of well-verified
falsehoods (the caloric theory's confirmed predictions) and of truths
unrecognised for centuries (heliocentrism before its confirmations).
That such conjunctions are \emph{coherent} shows that the success
condition of $R$ and the obtaining of $V(R)$ are distinct conditions;
identifying them would render the conjunctions contradictory, which they
are not. Moreover the proposed reduction is unstable: to state $V(R)$ at
all is to use language descriptively, hence (by \ref{C1} applied to
$V(R)$) to presuppose a non-epistemic success condition for $V(R)$
itself. The anti-realist reduction therefore either fails to capture the
ordinary distinction between being right and being justified, or
regresses, presupposing at the meta-level exactly the realist success
condition it sought to eliminate at the object level.
$\qquad\blacksquare$

Lemma~3 secures \ref{P2}: since the success condition cannot be reduced to
our epistemic situation, it is fixed by how things stand independently of
that situation. With \ref{P2} in hand, \ref{C1} follows.

\subsection{What has and has not been shown}

It is essential to state the verdict on the first sub-question precisely,
to avoid overclaiming.

\begin{description}[leftmargin=1.6em,style=nextline]
\item[Verdict on realism.] The descriptive use of language to talk about
the world \emph{does} supply an argument for realism. The argument is
\emph{conclusive for semantic realism understood as a presupposition}:
one cannot coherently use language descriptively while denying that
descriptions answer to a mind-independent standard of correctness
(\ref{C1}, Lemma~3).
\end{description}

Two honest qualifications. First, the argument is \emph{transcendental},
not \emph{metaphysically demonstrative}: it shows that realism is a
presupposition of a practice (description) in which we are inescapably
engaged, not that realism could be established from a standpoint outside
all description. An opponent willing to abandon the descriptive use of
language entirely would be untouched---but such an opponent could not
state, defend, or even understand any thesis, including anti-realism, for
to state a thesis is to describe. The presupposition is thus
\emph{practically inescapable}, which is the strongest status a thesis at
this level can have. Second, the argument secures ontological realism
(thesis~(a)) only as far as semantic realism (thesis~(b)) requires it:
there must be \emph{some} mind-independent matters that fix success
conditions. It does not by itself adjudicate the reality of any
\emph{particular} posited entity (electrons, fields). That finer question
is settled, when it is settled at all, by the abductive considerations of
\S4, not by the present transcendental argument. Keeping these two
results separate is what prevents the essay from sliding from a strong,
narrow conclusion to a weak, overreaching one.

\section{From realism to theoreticism}

We now have: descriptive language presupposes that descriptions answer to
a mind-independent world (semantic realism, with the minimal ontological
realism it entails). I argue that this result tells decisively against
operationalism and instrumentalism and in favour of theoreticism. The
argument has two parts: a negative part (the rejected doctrines are
untenable or self-undermining given realism) and a positive part
(theoreticism is the position that accommodates the realist datum).

\subsection{Against operationalism}

\begin{description}[leftmargin=1.6em,style=nextline]
\item[Claim 4.] Operationalism is false: scientific concepts are not
exhausted by the operations used to measure them.
\end{description}

\noindent\emph{Argument.} (i)~\emph{The over-multiplication objection.}
If a concept just is its measuring operations, then distinct operations
define distinct concepts. Length-by-rod and length-by-triangulation and
length-by-light-travel-time would be three concepts, not one quantity
measured three ways. But science treats them as measuring \emph{the same}
length, and asks whether the methods \emph{agree}---a question that
presupposes a single quantity, fixed independently of any one operation,
that the operations are answerable to. By Lemma~2 the report ``these
operations agree'' is bipolar: the operations \emph{might disagree}, and
when they do we conclude that one is in error, not that the world changed
its number of lengths. This is intelligible only if length is a
mind-independent magnitude (semantic realism, \ref{C1}) that no single
operation defines. (ii)~\emph{The operation-presupposes-theory
objection.} Specifying a measuring operation already deploys descriptive,
theory-laden concepts (``rigid rod'', ``simultaneous'', ``undisturbed
sample''); the operation cannot define the concept because understanding
the operation already requires the concept. Hence operationalism either
fails to individuate concepts correctly (objection~i) or presupposes the
very descriptive/theoretical content it claims to eliminate
(objection~ii). In both cases it founders on the realist datum of \S3.
$\qquad\blacksquare$

Note the dependence on the earlier results: the cogency of ``the
operations might disagree, and then one is wrong'' is exactly the
bipolarity of Lemma~2 plus the mind-independence of \ref{C1}. Without
those, the operationalist's identification of concept with operation
could not be faulted.

\subsection{Against instrumentalism}

\begin{description}[leftmargin=1.6em,style=nextline]
\item[Claim 5.] Instrumentalism is untenable: theories cannot be
construed as mere predictive instruments with no truth-value.
\end{description}

\noindent\emph{Argument.} (i)~\emph{Observation statements are already
descriptive.} The predictions an instrumentalist grants to theories are
\emph{statements about observables}---``the pointer reads $5$'', ``the
line appears at $\lambda$''. By Lemmas~1--2 these are descriptions:
bipolar, truth-evaluable, answerable to a mind-independent state of
affairs (\ref{C1}). So truth is already in play at the observational
level the instrumentalist accepts; the instrumentalist cannot quarantine
truth to ``observations'' while denying it to ``theory'', because both are
made of the same descriptive fabric. (ii)~\emph{No principled
theory/observation line.} What counts as ``observable'' is itself fixed by
theory (whether a cloud-chamber track, a Geiger count, or a stellar
parallax counts as observing depends on accepted theory); there is no
theory-neutral observational stratum to which truth could be confined and
beyond which it could be withheld. (iii)~\emph{Instruments are not
asserted, but theories are.} A genuine instrument (a slide rule) is used,
not believed; it is neither true nor false. But scientists \emph{assert}
theories, deduce consequences from them, regard them as
\emph{contradicting} rival theories, and revise them under refutation. The
argumentative function \ref{F4}---inference, contradiction,
refutation---operates on theories exactly as on any body of descriptions,
and \ref{F4} presupposes \ref{F3} (\S2.1). A device that can be
\emph{contradicted by evidence} is being treated as a description, not as
a mere instrument. Hence the instrumentalist's own practice (testing,
refuting, preferring) is unintelligible unless theories are
truth-apt.$\qquad\blacksquare$

A residual instrumentalist move must be blocked: ``Grant that predictions
are true or false, but theories are merely \emph{efficient summaries} of
such predictions, with no surplus content to be true \emph{of}.'' This
fails on two counts. First, a theory that genuinely had no content beyond
its actual predictions could not be \emph{tested by new} predictions, yet
testability by novel cases is the mark of a theory; the surplus content is
what does the predicting. Second, by \ref{C1} the surplus descriptive
content, being descriptive, has its own bipolar success condition: it is
true or false of the world. One cannot enjoy the explanatory and
predictive fertility that comes \emph{from} treating a theory as
describing hidden structure while disowning the truth-aptness that such
describing carries.

\subsection{For theoreticism}

\begin{description}[leftmargin=1.6em,style=nextline]
\item[Claim 6.] Theoreticism is the position that accommodates the
results of \S2--\S3 and survives the objections that sink operationalism
and instrumentalism.
\end{description}

\noindent\emph{Argument.} Theoreticism asserts two things: (T1)~inquiry is
conducted from within a framework of theories (there is no
theory-neutral, purely operational or purely observational ground floor);
and (T2)~the aim of theories is genuine explanation, i.e.\ true
description of the structures and mechanisms that produce the phenomena,
not mere correlation of observables.

\emph{(T1) is established by the failures of the rivals.} The
operation-presupposes-theory objection (\S4.1.ii) and the
no-theory-neutral-observation objection (\S4.2.ii) are exactly the
positive content of (T1): measurement and observation are theory-laden,
so theory is prior to, not reducible to, operation and observation.

\emph{(T2) is established by the descriptive/realist datum.} By \ref{C1},
the theoretical statements scientists use are descriptions answering to a
mind-independent world; their success condition is that the world be as
they say. To say the world is thus-and-so at the level of unobserved
structure \emph{is} to explain why the observable correlations hold
(they hold because that structure obtains). The instrumentalist aim of
``mere correlation'' is therefore strictly weaker than, and parasitic on,
the explanatory aim: the correlations the instrumentalist prizes are the
\emph{observable shadow} of the structure the theoretician describes, and
are predicted \emph{because} that structure is posited (\S4.2, residual
move). Thus the descriptive function, once admitted, already commits
inquiry to the explanatory, structure-describing aim that defines
theoreticism.$\qquad\blacksquare$

\subsection{The strength of the endorsement, stated honestly}

The negative results (Claims~4 and~5) are as strong as the realist datum
that drives them: given that descriptive language presupposes realism
(\S3, conclusive), operationalism and instrumentalism are refuted, not
merely disfavoured, in the strong forms defined in \S1. (T1) is likewise
established outright. The one component that is \emph{abductive} rather
than demonstrative is the inference from ``theoretical statements are
true-or-false descriptions of mind-independent structure'' to ``their
characteristic aim is, and should be, \emph{explanation by true
structural description}''. That theories so understood \emph{can}
explain, and that the explanatory construal best accounts for the actual
practices of testing, unification, novel prediction, and inter-method
agreement, is an inference to the best explanation. It is strong---no
rival construal accounts for those practices as well---but it is not a
deduction, and I do not present it as one. Accordingly:

\begin{description}[leftmargin=1.6em,style=nextline]
\item[Final verdict.] The descriptive use of language to talk about the
world \emph{does} supply an argument in favour of realism (conclusive for
semantic realism as a practically inescapable presupposition,
\S3). This realism in turn \emph{does} justify replacing operationalism
and instrumentalism by theoreticism: the realist datum \emph{refutes}
operationalism and instrumentalism (Claims~4--5) and \emph{establishes}
the theory-ladenness limb (T1) of theoreticism; and it \emph{strongly
supports}, by inference to the best explanation, the explanatory-aim limb
(T2). The replacement is therefore justified---deductively as regards the
rejection of the rivals and the priority of theory, and abductively (at
the highest available grade) as regards the explanatory aim.
\end{description}

\section{Objections and replies}

\paragraph{Objection 1 (constructivism).} ``The world we describe is
constituted by our conceptual scheme; there is no scheme-independent
success condition, so \ref{P2} fails.'' \emph{Reply.} The claim ``the
world is constituted by our scheme'' is itself a description (it says how
things stand with world and scheme) and so, by Lemma~2, is bipolar: it
could be false. But its content---that there is \emph{no}
scheme-independent fact---is precisely the denial of a scheme-independent
success condition for itself, which makes its own bipolarity
unintelligible. The position is self-refuting in the way diagnosed in
Lemma~3: stating it presupposes what it denies. At most it relocates
mind-independence (the scheme, and the fact of its constituting, become
the mind-independent matters), which still vindicates \ref{P2}.

\paragraph{Objection 2 (van Fraassen's constructive empiricism).} ``One
may accept that theories are truth-apt (so instrumentalism in the strict
sense is false) yet hold that the aim of science is only empirical
adequacy, not truth; this blocks (T2) without operationalist or
instrumentalist commitments.'' \emph{Reply.} This concedes the negative
results of \S4 and the whole of \S3; it disputes only the abductive limb
(T2). I have already classified (T2) as abductive, so the objection does
not touch the demonstrative core. Against (T2) specifically: empirical
adequacy is defined as truth about \emph{observables}, and \S4.2.ii
showed there is no theory-neutral fixing of ``observable''; so ``aim only
at empirical adequacy'' has no stable target distinct from ``aim at
truth'' once the observable/unobservable line is seen to be drawn by
theory. The best explanation of the actual selection pressures in science
(novel prediction, unification, inter-method agreement, the abandonment
of empirically adequate-so-far theories under deeper theoretical
criticism) remains the explanatory-truth aim of (T2). The objection
lowers, but does not remove, the force of the (already explicitly
abductive) endorsement of (T2).

\paragraph{Objection 3 (deflationism about truth).} ``If truth is merely
disquotational, the appeal to truth-aptness proves nothing
metaphysical.'' \emph{Reply.} The argument of \S3 does not require an
inflationary theory of truth; it requires only the bipolarity of
description (Lemma~2) and the non-identity of success and verification
(Lemma~3). Even on a deflationary reading, ``$R$ is true'' is
intersubstitutable with $R$ and so inherits $R$'s bipolar,
verification-transcendent success condition; the disquotational schema
``$R$ is true iff $R$'' is exactly the mind-independence of the success
condition rewritten. Deflationism therefore leaves \ref{C1} standing.

\section{Conclusion}

The descriptive function of language is the use of declarative sentences
to say how things stand, and it is the seat of truth, error, and
inference (\S2). Built into that function is a standard of correctness
that the act of describing cannot itself secure (Lemmas~2--3); to use
language descriptively is therefore to presuppose a mind-independent
world to which descriptions answer (\S3). This is a genuine and, at the
level of presupposition, conclusive argument for realism. The realist
datum refutes operationalism (which cannot individuate or even specify its
concepts without the descriptive content it denies) and instrumentalism
(whose own predictions, tests, and refutations are descriptive and
truth-apt), and it establishes that inquiry is irreducibly theory-laden.
It strongly supports, by inference to the best explanation, that the aim
of theory is genuine explanation by true structural description. Hence the
replacement of operationalism and instrumentalism by theoreticism is
justified: deductively in its negative and theory-ladenness components,
and abductively at the highest available grade in its explanatory-aim
component.

\end{document}
